% !TEX program = xelatex
% ============================================================
%  重庆财经学院本科毕业论文 LaTeX 模板
%  依据：《重庆财经学院毕业论文（设计）指导手册》
%  编译方式：XeLaTeX（需编译两次以生成目录和交叉引用）
% ============================================================

\documentclass[12pt,a4paper]{ctexart}

% ---------- 宏包 ----------
\usepackage{geometry}
\usepackage{setspace}
\usepackage{fancyhdr}
\usepackage{caption}
\usepackage{tocloft}
\usepackage{enumitem}
\usepackage{amsmath,amssymb}
\usepackage{graphicx}
\usepackage{booktabs}
\usepackage{hanging}
\usepackage{hyperref}
\hypersetup{colorlinks=true,linkcolor=black,citecolor=black,urlcolor=black}

% ---------- 页面设置 ----------
% 页边距：左2.8cm，右/上/下各2.5cm；页眉1.5cm，页脚1.75cm
\geometry{
  a4paper,
  left=2.8cm,
  right=2.5cm,
  top=2.5cm,
  bottom=2.5cm,
  headheight=1.5cm,
  footskip=1.75cm,
}

% ---------- 行距：固定值26磅 ----------
% Word固定值26磅 ≈ \baselineskip=26pt
\AtBeginDocument{\setlength{\baselineskip}{26pt}}
% 表格内使用单倍行距
\renewcommand{\arraystretch}{1.0}

% ---------- 字体设置 ----------
% ctex 自动选择中文字体：\songti(宋体) \heiti(黑体)
% 英文默认使用 Latin Modern；如需 Times New Roman，取消下一行注释（需系统已安装该字体）
% \setmainfont{Times New Roman}

% ---------- 页眉页脚 ----------
\pagestyle{fancy}
\fancyhf{}
\fancyfoot[C]{\thepage}
\renewcommand{\headrulewidth}{0pt}
\renewcommand{\footrulewidth}{0pt}

% ============================================================
%  自定义四级标题命令
%  一级：一、二、三…  黑体小三号加粗居中
%  二级：（一）（二）… 黑体四号首行缩进2字符
%  三级：1. 2. 3…    黑体小四号首行缩进2字符
%  四级：（1）（2）…  宋体小四号加粗首行缩进2字符
% ============================================================

% 章节阿拉伯数字计数器（用于图表和公式编号）
\newcounter{chapternum}

% 一级标题
\newcounter{cnt1}
\renewcommand{\thecnt1}{\chinese{cnt1}}
\newcommand{\h1}[1]{%
  \refstepcounter{cnt1}%
  \stepcounter{chapternum}%
  \par\vspace{1em}%
  {\centering\heiti\zihao{-3}\bfseries \thecnt1、#1\par}%
  \vspace{0.5em}%
  \addcontentsline{toc}{section}{\protect\numberline{\thecnt1、}#1}%
  \setcounter{cnt2}{0}\setcounter{cnt3}{0}\setcounter{cnt4}{0}%
  \setcounter{figure}{0}\setcounter{table}{0}\setcounter{equation}{0}%
}

% 二级标题
\newcounter{cnt2}
\renewcommand{\thecnt2}{（\chinese{cnt2}）}
\newcommand{\h2}[1]{%
  \refstepcounter{cnt2}%
  \par\vspace{0.5em}%
  {\par\setlength{\parindent}{2em}\heiti\zihao{4}\thecnt2#1\par}%
  \vspace{0.3em}%
  \addcontentsline{toc}{subsection}{\protect\numberline{\thecnt2}#1}%
  \setcounter{cnt3}{0}\setcounter{cnt4}{0}%
}

% 三级标题
\newcounter{cnt3}
\renewcommand{\thecnt3}{\arabic{cnt3}.}
\newcommand{\h3}[1]{%
  \refstepcounter{cnt3}%
  \par\vspace{0.3em}%
  {\par\setlength{\parindent}{2em}\heiti\zihao{-4}\thecnt3#1\par}%
  \setcounter{cnt4}{0}%
}

% 四级标题
\newcounter{cnt4}
\renewcommand{\thecnt4}{（\arabic{cnt4}）}
\newcommand{\h4}[1]{%
  \refstepcounter{cnt4}%
  \par\vspace{0.2em}%
  {\par\setlength{\parindent}{2em}\songti\zihao{-4}\bfseries\thecnt4#1\par}%
}

% ============================================================
%  图表设置：图 1-1 / 表 1-1，按章节编号
% ============================================================
\renewcommand{\thefigure}{\arabic{chapternum}-\arabic{figure}}
\renewcommand{\thetable}{\arabic{chapternum}-\arabic{table}}
\captionsetup{
  font={heiti,zihao={5}},
  labelsep=space,
  justification=centering,
  singlelinecheck=false,
}
% 表题在上、图题在下的使用约定：
%   表格：\begin{table} \caption{...} \begin{tabular}... \end{tabular} \end{table}
%   图片：\begin{figure} \includegraphics{...} \caption{...} \end{figure}

% ============================================================
%  公式设置：编号（1.1）右对齐，按章节编号
% ============================================================
\numberwithin{equation}{chapternum}
\renewcommand{\theequation}{\arabic{chapternum}.\arabic{equation}}

% ============================================================
%  目录设置：只显示到二级，点号前导符
% ============================================================
\setcounter{tocdepth}{2}
\renewcommand{\cftsecfont}{\songti\zihao{-4}}
\renewcommand{\cftsubsecfont}{\songti\zihao{-4}}
\renewcommand{\cftsecpagefont}{\songti\zihao{-4}}
\renewcommand{\cftsubsecpagefont}{\songti\zihao{-4}}
\setlength{\cftsubsecindent}{2em}
\renewcommand{\cftsecdotsep}{\cftdotsep}
\renewcommand{\cftsubsecdotsep}{\cftdotsep}
\renewcommand{\contentsname}{目\hspace{2em}录}

% ============================================================
%  正文段落：宋体小四号，首行缩进2字符，两端对齐
% ============================================================
\setlength{\parindent}{2em}
\renewcommand{\CJKglue}{\hskip 0pt plus 0.08\baselineskip}

% ============================================================
%  参考文献环境：悬挂缩进2字符，五号字
% ============================================================
\newenvironment{therefs}{%
  \begin{enumerate}[
    label={[\arabic*]},
    font=\songti\zihao{5},
    leftmargin=2em,
    itemindent=-2em,
    labelsep=0.5em,
    itemsep=0pt,
    parsep=0pt,
    topsep=0pt,
  ]
}{%
  \end{enumerate}
}

% ============================================================
%  文档开始
% ============================================================
\begin{document}

% ============================================================
%  封面（无页码）
%  格式：宋体/Times New Roman，四号，加粗，居中
% ============================================================
\begin{titlepage}
\thispagestyle{empty}
\begin{center}
\vspace*{3cm}
{\songti\zihao{4}\bfseries 重庆财经学院}\\[2cm]
{\songti\zihao{4}\bfseries 本科毕业论文（设计）}\\[3cm]
% 论文题目
{\songti\zihao{4}\bfseries 论文题目：}\\[0.5cm]
{\songti\zihao{4}\bfseries （此处填写论文题目）}\\[3cm]
% 信息栏
\begin{tabular}{c}
{\songti\zihao{4}\bfseries 学　　院：\underline{\hspace{6em}}}\\[0.8em]
{\songti\zihao{4}\bfseries 年　　级：\underline{\hspace{6em}}}\\[0.8em]
{\songti\zihao{4}\bfseries 专　　业：\underline{\hspace{6em}}}\\[0.8em]
{\songti\zihao{4}\bfseries 学　　号：\underline{\hspace{6em}}}\\[0.8em]
{\songti\zihao{4}\bfseries 姓　　名：\underline{\hspace{6em}}}\\[0.8em]
{\songti\zihao{4}\bfseries 指导教师：\underline{\hspace{6em}}}\\[0.8em]
\end{tabular}
\vfill
{\songti\zihao{4}\bfseries 二〇二六年六月}
\end{center}
\end{titlepage}

% ============================================================
%  前置部分：大写罗马数字页码
% ============================================================
\pagenumbering{Roman}

% ---------- 中文摘要 ----------
\newpage
\begin{center}
{\heiti\zihao{-3}\bfseries 摘\hspace{2em}要}
\end{center}
\vspace{1em}
{\songti\zihao{-4}
（此处填写中文摘要内容。首行缩进2字符，固定行距26磅。）
}\par
\vspace{1em}
\noindent{\heiti\zihao{-4}关键词：}{\songti\zihao{-4}关键词1；关键词2；关键词3；关键词4}

% ---------- 英文摘要 ----------
\newpage
\begin{center}
{\zihao{-3}\bfseries Abstract}
\end{center}
\vspace{1em}
{\zihao{-4}
（Fill in the English abstract here. First line indent 2 characters, fixed line spacing 26pt.）
}\par
\vspace{1em}
\noindent{\zihao{-4}\bfseries Keywords: }{\zihao{-4}keyword1; keyword2; keyword3; keyword4}

% ---------- 目录 ----------
\newpage
\tableofcontents

% ============================================================
%  正文部分：阿拉伯数字页码，从1开始
% ============================================================
\newpage
\pagenumbering{arabic}

% ---------- 第一章 ----------
\h1{绪论}

\h2{研究背景}

{\songti\zihao{-4}
（此处填写正文内容。宋体小四号，首行缩进2字符，两端对齐，固定行距26磅。正文引用参考文献使用上标格式，如\cite{ref1}、\cite{ref2,ref3}。）
}\par

\h2{研究意义}

\h3{理论意义}

\h3{现实意义}

\h4{第一点}

% ---------- 第二章 ----------
\h1{文献综述}

% ---------- 第三章 ----------
\h1{研究设计与方法}

% 图表示例：图题在图下方
\begin{figure}[htbp]
\centering
\includegraphics[width=0.8\textwidth]{example-image}
\caption{图题示例}
\end{figure}

% 图表示例：表题在表上方，推荐三线表
\begin{table}[htbp]
\centering
\caption{表题示例}
\begin{tabular}{ccc}
\toprule
列1 & 列2 & 列3 \\
\midrule
数据1 & 数据2 & 数据3 \\
数据4 & 数据5 & 数据6 \\
\bottomrule
\end{tabular}
\end{table}

% 公式示例：编号（3.1）右对齐
\begin{equation}
E = mc^2
\end{equation}

% ---------- 第四章 ----------
\h1{实证分析}

% ---------- 第五章 ----------
\h1{结论与展望}

% ============================================================
%  参考文献（另起一页，悬挂缩进2字符，五号字）
% ============================================================
\newpage
\begin{center}
{\heiti\zihao{-3}\bfseries 参考文献}
\end{center}
\vspace{1em}
\begin{therefs}
\item 作者. 文献标题[J]. 期刊名, 年份, 卷(期): 起止页码.
\item 作者. 书名[M]. 出版地: 出版社, 年份.
\item 作者. 论文题目[D]. 保存地: 保存单位, 年份.
\end{therefs}

% ============================================================
%  附录（如有）
% ============================================================
\newpage
\begin{center}
{\heiti\zihao{-3}\bfseries 附录\hspace{1em}A}
\end{center}
\vspace{0.5em}
\begin{center}
{\heiti\zihao{4}附录标题}
\end{center}
\vspace{0.5em}
{\songti\zihao{-4}
（附录条文内容，宋体小四号，首行缩进2字符。）
}\par

% ============================================================
%  文档结束
% ============================================================
\end{document}
